\section{Assign roles inside the leaf function}\label{vm-registers}
The public calling convention describes values at entry and return. Inside the guided
function, saved registers hold a private protocol. The prologue saves their original
full values before assigning these roles.

\begin{table}[htbp]
  \centering
  \caption{Register roles during the guided call. The bitmap uses all sixty-four register bits.}
  \tagpdfsetup{table/header-rows={1}}
  \begin{tabular}{ll}
    \toprule
    Registers                & Meaning                                        \\
    \midrule
    \texttt{s0}, \texttt{s1} & Full code base and byte length                 \\
    \texttt{s2}, \texttt{s3} & Full descriptor and operand-stack addresses    \\
    \texttt{s4}, \texttt{s5} & Capacity in words and instruction-start bitmap \\
    \texttt{s6}, \texttt{s7} & Phase cursor or guest PC, and live depth       \\
    \texttt{v0}, \texttt{v1} & Status and retired guest count                 \\
    \texttt{a2}, \texttt{a3} & Fault offset and remaining guest budget        \\
    \bottomrule
  \end{tabular}
\end{table}

The supplied \texttt{targets\_begin} and \texttt{execute\_begin} entries each reset
\texttt{s6} to zero. During validation it is a decode cursor; after the second reset
it becomes guest PC. No validation record increments \texttt{v1}. The fault register
begins at minus one and changes only when a failure is selected.

Every unfinished region initially branches to \texttt{vm\_commit}. Replace that
branch and its \texttt{nop}, leaving the named boundary comments and surrounding
labels. The resulting partial call can return cleanly before its full semantics are
implemented. Intermediate observations below stop at \texttt{vm\_commit} before its
stores and restores. At that point saved registers still contain interpreter roles.
After returning to \texttt{vm\_ready}, they contain the caller's restored values.
Do not confuse a partial decode cursor with a completed public output PC.

The guided code includes a \texttt{nop} after each retained delayed transfer. A
future native instruction placed in one of those slots would execute according to
that MIPS instruction's delay behavior. Preserve the listed slots while developing
the guest semantics; the bytecode language has no guest delay slots.
